\section{Conclusion}
\label{sec:conclusion}

We asked whether Vintern-1B can be improved on AutoViVQA by training only a small
fraction of its parameters, and where the bottleneck lies. A frozen InternViT-300M
and frozen Qwen2-0.5B, connected by a $0.78\%$ mean-pooled bridge at one tile and
adapted with a $0.23\%$ rank-16 attention LoRA on the decoder, matches or beats
the fully fine-tuned Vintern-1B on every generation metric and beats ViMoE-VQA on
corpus CIDEr-D / BLEU-4 / ROUGE-L --- at ${\sim}1\%$ trainable parameters, one
16\,GB GPU, no backbone fine-tuning, and with held-out test numbers matching
validation.

Across a six-question ablation ladder, four vision- and training-side levers ---
a $10\times$-larger bridge, more tiles, a per-question routing policy,
multi-reference training, projector-level alignment --- produce no F1 improvement
(alignment is an exact null). The only lever that helps is a small LoRA on the
frozen decoder; it helps on every bridge, works only through the attention
projections (feed-forward LoRA destabilises training), and collapses the five
bridges into a $0.6$-point F1 band. For this VLM class, the ceiling on
token-level phrasing is the frozen decoder's attention capacity, not the visual
pipeline --- and on the same benchmark, question type carries no signal about the
optimal visual-compute action, a direct counterpoint to ``reasoning-aware''
routing accounts.

We release the code, the leak-free grouped split, the per-question oracle tables,
and all checkpoints. The natural next step is to vary the one component we kept
frozen and small --- the decoder --- and test whether a larger frozen Vietnamese
decoder closes the remaining token-F1 gap without any of the levers that failed
here.
